% !TeX spellcheck = pl_PL
%%%%%%%%%%%%%%%%%%%%%%%%%%%%%%%%%%%%%%%%%%%
%                                        %
% Szablon pracy dyplomowej inzynierskiej %
% zgodny  z aktualnymi  przepisami  SZJK %
%                                        %
%%%%%%%%%%%%%%%%%%%%%%%%%%%%%%%%%%%%%%%%%%
%                                        %
%  (c) Krzysztof Simiński, 2018-2023     %
%                                        %
%%%%%%%%%%%%%%%%%%%%%%%%%%%%%%%%%%%%%%%%%%
%                                        %
% Najnowsza wersja szablonów jest        %
% podstępna pod adresem                  %
% github.com/ksiminski/polsl-aei-theses  %
%                                        %
%%%%%%%%%%%%%%%%%%%%%%%%%%%%%%%%%%%%%%%%%%
%
%
% Projekt LaTeXowy zapewnia odpowiednie formatowanie pracy,
% zgodnie z wymaganiami Systemu zapewniania jakości kształcenia.
% Proszę nie zmieniać ustawień formatowania (np. fontu,
% marginesów, wytłuszczeń, kursywy itd. ).
%
% Projekt można kompilować na kilka sposobów.
%
% 1. kompilacja pdfLaTeX
%
% pdflatex main
% bibtex   main
% pdflatex main
% pdflatex main
%
%
% 2. kompilacja XeLaTeX
%
% Kompilatacja przy użyciu XeLaTeXa różni się tym, że na stronie
% tytułowej używany jest font Calibri. Wymaga to jego uprzedniego
% zainstalowania.
%
% xelatex main
% bibtex  main
% xelatex main
% xelatex main
%
%
%%%%%%%%%%%%%%%%%%%%%%%%%%%%%%%%%%%%%%%%%%%%%%%%%%%%%
% W przypadku pytań, uwag, proszę pisać na adres:   %
%      krzysztof.siminski(małpa)polsl.pl            %
%%%%%%%%%%%%%%%%%%%%%%%%%%%%%%%%%%%%%%%%%%%%%%%%%%%%%
%
% Chcemy ulepszać szablony LaTeXowe prac dyplomowych.
% Wypełniając ankietę spod poniższego adresu pomogą
% Państwo nam to zrobić. Ankieta jest całkowicie
% anonimowa. Dziękujemy!


% https://docs.google.com/forms/d/e/1FAIpQLScyllVxNKzKFHfILDfdbwC-jvT8YL0RSTFs-s27UGw9CKn-fQ/viewform?usp=sf_link
%
%%%%%%%%%%%%%%%%%%%%%%%%%%%%%%%%%%%%%%%%%%%%%%%%%%%%%%%%%%%%%%%%%%%%%%%%%

%%%%%%%%%%%%%%%%%%%%%%%%%%%%%%%%%%%%%%%%%%%%%%%
%                                             %
% PERSONALIZACJA PRACY – DANE PRACY           %
%                                             %
%%%%%%%%%%%%%%%%%%%%%%%%%%%%%%%%%%%%%%%%%%%%%%%

% Proszę wpisać swoje dane w poniższych definicjach.

% TODO
% dane autora
\newcommand{\FirstNameAuthor}{Michał}
\newcommand{\SurnameAuthor}{Jagoda}
\newcommand{\IdAuthor}{300405}   % numer albumu  (bez $\langle$ i $\rangle$)

% drugi autor:
%\newcommand{\FirstNameCoauthor}{Imię}   % Jeżeli jest drugi autor, to tutaj należy podać imię.
%\newcommand{\SurnameCoauthor}{Nazwisko} % Jeżeli jest drugi autor, to tutaj należy podać nazwisko.
%\newcommand{\IdCoauthor}{$\langle$wpisać właściwy$\rangle$}  % numer albumu drugiego autora (bez $\langle$ i $\rangle$)
% Gdy nie ma drugiego autora, należy zostawić poniższe definicje puste, jak poniżej. Gdy jest drugi autor, należy zakomentować te linie.
\newcommand{\FirstNameCoauthor}{} % Jeżeli praca ma tylko jednego autora, to dane drugiego autora zostają puste.
\newcommand{\SurnameCoauthor}{}   % Jeżeli praca ma tylko jednego autora, to dane drugiego autora zostają puste.
\newcommand{\IdCoauthor}{}  % Jeżeli praca ma tylko jednego autora, to dane drugiego autora zostają puste.
%%%%%%%%%%

\newcommand{\Supervisor}{Dr hab. inż. Karolina Nurzyńska}     % dane promotora (bez $\langle$ i $\rangle$)
\newcommand{\Title}{Projekt i implementacja aplikacji webowej do transkrypcji nagrań fortepianowych na nuty w formacie PDF z wykorzystaniem zaawansowanych modeli AI}           % tytuł pracy po polsku
\newcommand{\TitleAlt}{Design and implementation of a web application for transcribing piano recordings into musical notes in PDF format using advanced AI models}                     % thesis title in English
\newcommand{\Program}{Informatyka}            % kierunek studiów  (bez $\langle$ i $\rangle$)
\newcommand{\Specialisation}{Grafika Komputerowa i Oprogramowanie}     % specjalność  (bez $\langle$ i $\rangle$)
\newcommand{\Departament}{Algorytmiki i Oprogramowania}        % katedra promotora  (bez $\langle$ i $\rangle$)

% Jeżeli został wyznaczony promotor pomocniczy lub opiekun, proszę go/ją wpisać ...
%\newcommand{\Consultant}{$\langle$stopień naukowy imię i nazwisko$\rangle$} % dane promotora pomocniczego, opiekuna (bez $\langle$ i $\rangle$)
% ... w przeciwnym razie proszę zostawić puste miejsce jak poniżej:
\newcommand{\Consultant}{} % brak promotowa pomocniczego / opiekuna

% koniec fragmentu do modyfikacji
%%%%%%%%%%%%%%%%%%%%%%%%%%%%%%%%%%%%%%%%%%


%%%%%%%%%%%%%%%%%%%%%%%%%%%%%%%%%%%%%%%%%%%%%%%
%                                             %
% KONIEC PERSONALIZACJI PRACY                 %
%                                             %
%%%%%%%%%%%%%%%%%%%%%%%%%%%%%%%%%%%%%%%%%%%%%%%

%%%%%%%%%%%%%%%%%%%%%%%%%%%%%%%%%%%%%%%%


%%%%%%%%%%%%%%%%%%%%%%%%%%%%%%%%%%%%%%%%%%%%%%%
%                                             %
% PROSZĘ NIE MODYFIKOWAĆ PONIŻSZYCH USTAWIEŃ! %
%                                             %
%%%%%%%%%%%%%%%%%%%%%%%%%%%%%%%%%%%%%%%%%%%%%%%



\documentclass[a4paper,twoside,12pt]{book}
\usepackage[utf8]{inputenc}                                      
\usepackage[T1]{fontenc}  
\usepackage{amsmath,amsfonts,amssymb,amsthm}
\usepackage[british,polish]{babel} 
\usepackage{indentfirst}
\usepackage{xurl}
\usepackage{xstring}
\usepackage{ifthen}



\usepackage{ifxetex}

\ifxetex
	\usepackage{fontspec}
	\defaultfontfeatures{Mapping=tex—text} % to support TeX conventions like ``——-''
	\usepackage{xunicode} % Unicode support for LaTeX character names (accents, European chars, etc)
	\usepackage{xltxtra} % Extra customizations for XeLaTeX
\else
	\usepackage{lmodern}
\fi



\usepackage[margin=2.5cm]{geometry}
\usepackage{graphicx} 
\usepackage{hyperref}
\usepackage{booktabs}
\usepackage{tikz}
\usepackage{pgfplots}
\usepackage{mathtools}
\usepackage{geometry}
\usepackage{subcaption}   % subfigures
\usepackage[page]{appendix} % toc,
\renewcommand{\appendixtocname}{Dodatki}
\renewcommand{\appendixpagename}{Dodatki}
\renewcommand{\appendixname}{Dodatek}

\usepackage{csquotes}
\usepackage[natbib=true,backend=bibtex,maxbibnames=99]{biblatex}  % kompilacja bibliografii BibTeXem
%\usepackage[natbib=true,backend=biber,maxbibnames=99]{biblatex}  % kompilacja bibliografii Biberem
\bibliography{biblio}

\usepackage{ifmtarg}   % empty commands  

\usepackage{setspace}
\onehalfspacing


\frenchspacing

%%%%%%%%%%%%%%%%%%%%%%%%%%%%%%%%%%
% środowiska dla definicji, twierdzenia, przykładu
\usepackage{amsthm}

\newtheorem{Definition}{Definicja}
\newtheorem{Example}{Przykład}
\newtheorem{Theorem}{Twierdzenie}
%%%%%%%%%%%%%%%%%%%%%%%%%%%%%%%%%%

%%%% TODO LIST GENERATOR %%%%%%%%%

\usepackage{color}
\definecolor{brickred}      {cmyk}{0   , 0.89, 0.94, 0.28}

\makeatletter \newcommand \kslistofremarks{\section*{Uwagi} \@starttoc{rks}}
  \newcommand\l@uwagas[2]
    {\par\noindent \textbf{#2:} %\parbox{10cm}
{#1}\par} \makeatother


\newcommand{\ksremark}[1]{%
{%\marginpar{\textdbend}
{\color{brickred}{[#1]}}}%
\addcontentsline{rks}{uwagas}{\protect{#1}}%
}

\newcommand{\comma}{\ksremark{przecinek}}
\newcommand{\nocomma}{\ksremark{bez przecinka}}
\newcommand{\styl}{\ksremark{styl}}
\newcommand{\ortografia}{\ksremark{ortografia}}
\newcommand{\fleksja}{\ksremark{fleksja}}
\newcommand{\pauza}{\ksremark{pauza `--', nie dywiz `-'}}
\newcommand{\kolokwializm}{\ksremark{kolokwializm}}
\newcommand{\cudzyslowy}{\ksremark{,,polskie cudzysłowy''}}

%%%%%%%%%%%%%% END OF TODO LIST GENERATOR %%%%%%%%%%%

\newcommand{\printCoauthor}{%		
    \StrLen{\FirstNameCoauthor}[\FNCoALen]
    \ifthenelse{\FNCoALen > 0}%
    {%
		{\large\bfseries\Coauthor\par}
	
		{\normalsize\bfseries \LeftId: \IdCoauthor\par}
    }%
    {}
} 

%%%%%%%%%%%%%%%%%%%%%
\newcommand{\autor}{%		
    \StrLen{\FirstNameCoauthor}[\FNCoALenXX]
    \ifthenelse{\FNCoALenXX > 0}%
    {\FirstNameAuthor\ \SurnameAuthor, \FirstNameCoauthor\ \SurnameCoauthor}%
	{\FirstNameAuthor\ \SurnameAuthor}%
}
%%%%%%%%%%%%%%%%%%%%%

\StrLen{\FirstNameCoauthor}[\FNCoALen]
\ifthenelse{\FNCoALen > 0}%
{%
\author{\FirstNameAuthor\ \SurnameAuthor, \FirstNameCoauthor\ \SurnameCoauthor}
}%
{%
\author{\FirstNameAuthor\ \SurnameAuthor}
}%

%%%%%%%%%%%% ZYWA PAGINA %%%%%%%%%%%%%%%
% brak kapitalizacji zywej paginy
\usepackage{fancyhdr}
\pagestyle{fancy}
\fancyhf{}
\fancyhead[LO]{\nouppercase{\it\rightmark}}
\fancyhead[RE]{\nouppercase{\it\leftmark}}
\fancyhead[LE,RO]{\it\thepage}


\fancypagestyle{tylkoNumeryStron}{%
   \fancyhf{} 
   \fancyhead[LE,RO]{\it\thepage}
}

\fancypagestyle{bezNumeracji}{%
   \fancyhf{} 
   \fancyhead[LE,RO]{}
}


\fancypagestyle{NumeryStronNazwyRozdzialow}{%
   \fancyhf{} 
   \fancyhead[LE]{\nouppercase{\autor}}
   \fancyhead[RO]{\nouppercase{\leftmark}} 
   \fancyfoot[CE, CO]{\thepage}
}


%%%%%%%%%%%%% OBCE WTRETY  
\newcommand{\obcy}[1]{\emph{#1}}
\newcommand{\english}[1]{{\selectlanguage{british}\obcy{#1}}}
%%%%%%%%%%%%%%%%%%%%%%%%%%%%%

% polskie oznaczenia funkcji matematycznych
\renewcommand{\tan}{\operatorname {tg}}
\renewcommand{\log}{\operatorname {lg}}

% jeszcze jakies drobiazgi

\newcounter{stronyPozaNumeracja}

%%%%%%%%%%%%%%%%%%%%%%%%%%% 
\newcommand{\printOpiekun}[1]{%		

    \StrLen{\Consultant}[\mystringlen]
    \ifthenelse{\mystringlen > 0}%
    {%
       {\large{\bfseries OPIEKUN, PROMOTOR POMOCNICZY}\par}
       
       {\large{\bfseries \Consultant}\par}
    }%
    {}
} 
%
%%%%%%%%%%%%%%%%%%%%%%%%%%%%%%%%%%%%%%%%%%%%%%
 
% Proszę nie modyfikować poniższych definicji!
\newcommand{\Author}{\FirstNameAuthor\ \MakeUppercase{\SurnameAuthor}} 
\newcommand{\Coauthor}{\FirstNameCoauthor\ \MakeUppercase{\SurnameCoauthor}}
\newcommand{\Type}{PROJEKT INŻYNIERSKI}
\newcommand{\Faculty}{Wydział Automatyki, Elektroniki i Informatyki} 
\newcommand{\Polsl}{Politechnika Śląska}
\newcommand{\Logo}{politechnika_sl_logo_bw_pion_pl.pdf}
\newcommand{\LeftId}{Nr albumu}
\newcommand{\LeftProgram}{Kierunek}
\newcommand{\LeftSpecialisation}{Specjalność}
\newcommand{\LeftSUPERVISOR}{PROWADZĄCY PRACĘ}
\newcommand{\LeftDEPARTMENT}{KATEDRA}
%%%%%%%%%%%%%%%%%%%%%%%%%%%%%%%%%%%%%%%%%%%%%%

%%%%%%%%%%%%%%%%%%%%%%%%%%%%%%%%%%%%%%%%%%%%%%%
%                                             %
% KONIEC USTAWIEŃ                             %
%                                             %
%%%%%%%%%%%%%%%%%%%%%%%%%%%%%%%%%%%%%%%%%%%%%%%




%%%%%%%%%%%%%%%%%%%%%%%%%%%%%%%%%%%%%%%%%%%%%%%
%                                             %
% MOJE PAKIETY, USTAWIENIA ITD                %
%                                             %
%%%%%%%%%%%%%%%%%%%%%%%%%%%%%%%%%%%%%%%%%%%%%%%

% Tutaj proszę umieszczać swoje pakiety, makra, ustawienia itd.


 
%%%%%%%%%%%%%%%%%%%%%%%%%%%%%%%%%%%%%%%%%%%%%%%%%%%%%%%%%%%%%%%%%%%%%
% listingi i fragmentu kodu źródłowego 
% pakiet: listings lub minted
% % % % % % % % % % % % % % % % % % % % % % % % % % % % % % % % % % % 

% biblioteka listings
\usepackage{listings}
\lstset{%
morekeywords={string,exception,std,vector},% słowa kluczowe rozpoznawane przez pakiet listings
language=C++,% C, Matlab, Python, SQL, TeX, XML, bash, ... – vide https://www.ctan.org/pkg/listings
commentstyle=\textit,%
identifierstyle=\textsf,%
keywordstyle=\sffamily\bfseries, %\texttt, %
%captionpos=b,%
tabsize=3,%
frame=lines,%
numbers=left,%
numberstyle=\tiny,%
numbersep=5pt,%
breaklines=true,%
escapeinside={@*}{*@},%
}

% % % % % % % % % % % % % % % % % % % % % % % % % % % % % % % % % % % 
% pakiet minted
%\usepackage{minted}

% pakiet wymaga specjalnego kompilowania:
% pdflatex -shell-escape main.tex
% xelatex  -shell-escape main.tex

%\usepackage[chapter]{minted} % [section]
%%\usemintedstyle{bw}   % czarno-białe kody 
%
%\setminted % https://ctan.org/pkg/minted
%{
%%fontsize=\normalsize,%\footnotesize,
%%captionpos=b,%
%tabsize=3,%
%frame=lines,%
%framesep=2mm,
%numbers=left,%
%numbersep=5pt,%
%breaklines=true,%
%escapeinside=@@,%
%}

%%%%%%%%%%%%%%%%%%%%%%%%%%%%%%%%%%%%%%%%%%%%%%%%%%%%%%%%%%%%%%%%%%%%%



%%%%%%%%%%%%%%%%%%%%%%%%%%%%%%%%%%%%%%%%%%%%%%%
%                                             %
% KONIEC MOICH USTAWIEŃ                       %
%                                             %
%%%%%%%%%%%%%%%%%%%%%%%%%%%%%%%%%%%%%%%%%%%%%%%



%%%%%%%%%%%%%%%%%%%%%%%%%%%%%%%%%%%%%%%%


\begin{document}
%\kslistofremarks

\frontmatter

%%%%%%%%%%%%%%%%%%%%%%%%%%%%%%%%%%%%%%%%%%%%%%%
%                                             %
% PROSZĘ NIE MODYFIKOWAĆ STRONY TYTUŁOWEJ!    %
%                                             %
%%%%%%%%%%%%%%%%%%%%%%%%%%%%%%%%%%%%%%%%%%%%%%%


%%%%%%%%%%%%%%%%%%  STRONA TYTUŁOWA %%%%%%%%%%%%%%%%%%%
\pagestyle{empty}
{
	\newgeometry{top=1.5cm,%
	             bottom=2.5cm,%
	             left=3cm,
	             right=2.5cm}
 
	\ifxetex 
	  \begingroup
	  \setsansfont{Calibri}
	   
	\fi 
	 \sffamily
	\begin{center}
	\includegraphics[width=50mm]{\Logo}
	 
	
	{\Large\bfseries\Type\par}
	
	\vfill  \vfill  
			 
	{\large\Title\par}
	
	\vfill  
		
	{\large\bfseries\Author\par}
	
	{\normalsize\bfseries \LeftId: \IdAuthor}

	\printCoauthor
	
	\vfill  		
 
	{\large{\bfseries \LeftProgram:} \Program\par} 
	
	{\large{\bfseries \LeftSpecialisation:} \Specialisation\par} 
	 		
	\vfill  \vfill 	\vfill 	\vfill 	\vfill 	\vfill 	\vfill  
	 
	{\large{\bfseries \LeftSUPERVISOR}\par}
	
	{\large{\bfseries \Supervisor}\par}
				
	{\large{\bfseries \LeftDEPARTMENT\ \Departament} \par}
		
	{\large{\bfseries \Faculty}\par}
		
	\vfill  \vfill  

    	
    \printOpiekun{\Consultant}
    
	\vfill  \vfill  
		
    {\large\bfseries  Gliwice \the\year}

   \end{center}	
       \ifxetex 
       	  \endgroup
       \fi
	\restoregeometry
}
  
%%%%%%%%%%%%%%%%%%%%%%%%%%%%%%%%%%%%%%%%%%%%%%%
%                                             %
% KONIEC STRONY TYTUŁOWEJ                     %
%                                             %
%%%%%%%%%%%%%%%%%%%%%%%%%%%%%%%%%%%%%%%%%%%%%%%  


\cleardoublepage

\rmfamily\normalfont
\pagestyle{empty}


%%% No to zaczynamy pisać pracę :-) %%%%

% TODO
\subsubsection*{Tytuł pracy} 
\Title

\subsubsection*{Streszczenie}  
(Streszczenie pracy – odpowiednie pole w systemie APD powinno zawierać kopię tego streszczenia.)

\subsubsection*{Słowa kluczowe} 
(2-5 slow (fraz) kluczowych, oddzielonych przecinkami)

\subsubsection*{Thesis title} 
\begin{otherlanguage}{british}
\TitleAlt
\end{otherlanguage}

\subsubsection*{Abstract} 
\begin{otherlanguage}{british}
(Thesis abstract – to be copied into an appropriate field during an electronic submission – in English.)
\end{otherlanguage}
\subsubsection*{Key words}  
\begin{otherlanguage}{british}
(2-5 keywords, separated by commas)
\end{otherlanguage}




%%%%%%%%%%%%%%%%%% SPIS TRESCI %%%%%%%%%%%%%%%%%%%%%%
% Add \thispagestyle{empty} to the toc file (main.toc), because \pagestyle{empty} doesn't work if the TOC has multiple pages
\addtocontents{toc}{\protect\thispagestyle{empty}}
\tableofcontents

%%%%%%%%%%%%%%%%%%%%%%%%%%%%%%%%%%%%%%%%%%%%%%%%%%%%%
\setcounter{stronyPozaNumeracja}{\value{page}}
\mainmatter
\pagestyle{empty}

\cleardoublepage

\pagestyle{NumeryStronNazwyRozdzialow}

%%%%%%%%%%%%%% wlasciwa tresc pracy %%%%%%%%%%%%%%%%%

% TODO
\chapter{Wstęp}
\label{ch:wstep}

\begin{itemize}
\item wprowadzenie w problem/zagadnienie
\item osadzenie problemu w dziedzinie
\item cel pracy
\item zakres pracy
\item zwięzła charakterystyka rozdziałów
\item jednoznaczne określenie wkładu autora, w przypadku prac wieloosobowych – tabela z autorstwem poszczególnych elementów pracy
\end{itemize}



% TODO
\chapter{Analiza tematu}

\section{Sformułowanie problemu}

Współczesna muzyka fortepianowa często występuje w formie cyfrowych nagrań audio, najczęściej w formacie MP3. Format ten, choć wygodny do odtwarzania, nie pozwala na łatwą edycję czy analizę struktury muzycznej. Muzycy, kompozytorzy, pedagodzy muzyczni oraz entuzjaści fortepianu potrzebują narzędzia, które umożliwi im konwersję plików audio do formatu MIDI, a następnie do czytelnych partytur w formacie PDF. 

Problem, którym zajmuje się niniejsza praca, polega na stworzeniu efektywnego, dokładnego i przyjaznego dla użytkownika systemu, który automatyzuje proces transkrypcji muzyki fortepianowej. System ten, nazwany TuneScript, ma za zadanie przekształcać nagrania fortepianowe z formatu MP3 do symbolicznej reprezentacji MIDI, a następnie generować profesjonalnie wyglądające partytury w formacie PDF.

Kluczowe wyzwania, które należy rozwiązać w ramach tego projektu, obejmują:

\begin{itemize}
    \item Precyzyjne rozpoznawanie wysokości dźwięków w nagraniach fortepianowych, z uwzględnieniem złożoności polifonii fortepianowej.
    \item Dokładne określanie czasu rozpoczęcia i zakończenia każdej nuty, co jest kluczowe dla zachowania rytmu i tempa utworu.
    \item Odwzorowanie dynamiki gry, która jest istotnym elementem ekspresji muzycznej na fortepianie.
    \item Efektywne przetwarzanie długich utworów muzycznych przy zachowaniu wysokiej jakości transkrypcji.
    \item Generowanie czytelnych i estetycznych partytur PDF, które będą użyteczne dla muzyków.
\end{itemize}

Rozwiązanie tego problemu ma potencjał znaczącego usprawnienia pracy muzyków, ułatwienia analizy muzycznej oraz wsparcia procesów edukacyjnych w dziedzinie muzyki fortepianowej.

\section{Osadzenie tematu w kontekście aktualnego stanu wiedzy}

Automatyczna transkrypcja muzyki (ATM, ang. \textit{Automatic Music Transcription}) jest jednym z najbardziej złożonych problemów w dziedzinie przetwarzania sygnałów muzycznych (MIR, ang. \textit{Music Information Retrieval}). W kontekście muzyki fortepianowej, ATM zyskuje szczególne znaczenie ze względu na popularność tego instrumentu oraz jego rolę w edukacji muzycznej i kompozycji.

Obecny stan wiedzy w dziedzinie ATM dla fortepianu charakteryzuje się następującymi aspektami:

\begin{enumerate}
    \item \textbf{Wykorzystanie głębokiego uczenia:} Współczesne podejścia do ATM fortepianowego opierają się głównie na zastosowaniu głębokich sieci neuronowych. Modele takie jak konwolucyjne sieci neuronowe (CNN) i rekurencyjne sieci neuronowe (RNN) wykazały znaczącą skuteczność w rozpoznawaniu struktur muzycznych w spektrogramach audio \cite{bib:hawthorne2018onsets}.
    
    \item \textbf{Modele hybrydowe:} Najnowsze badania skupiają się na modelach hybrydowych, łączących różne architektury sieci neuronowych. Na przykład, model "Onsets and Frames" \cite{bib:hawthorne2018onsets} wykorzystuje osobne sieci do detekcji początków dźwięków i ich trwania, co znacząco poprawiło dokładność transkrypcji.
    
    \item \textbf{Przetwarzanie w czasie rzeczywistym:} Rozwój technik przetwarzania sygnałów i optymalizacji algorytmów pozwala na coraz szybszą transkrypcję, zbliżającą się do przetwarzania w czasie rzeczywistym \cite{bib:kelz2016potential}.
    
    \item \textbf{Modelowanie akustyki fortepianu:} Zaawansowane modele uwzględniają specyfikę akustyki fortepianu, w tym zjawiska takie jak rezonans strun i wpływ pedałów na brzmienie \cite{bib:cheng2016attack}.
    
    \item \textbf{Transkrypcja ekspresji muzycznej:} Nowsze badania skupiają się nie tylko na rozpoznawaniu nut, ale także na odwzorowaniu dynamiki, artykulacji i innych aspektów ekspresji muzycznej \cite{bib:jeong2017deep}.
\end{enumerate}

Mimo znaczących postępów w dziedzinie ATM dla fortepianu, wciąż istnieją pewne ograniczenia i wyzwania:

\begin{itemize}
    \item \textbf{Dokładność w muzyce polifonicznej:} Transkrypcja złożonych struktur polifonicznych, charakterystycznych dla zaawansowanej literatury fortepianowej, wciąż stanowi wyzwanie dla obecnych systemów \cite{bib:benetos2013automatic}.
    
    \item \textbf{Rozpoznawanie bardzo cichych dźwięków:} Detekcja nut granych bardzo cicho, szczególnie w kontekście głośniejszego akompaniamentu, pozostaje problematyczna \cite{bib:sigtia2016end}.
    
    \item \textbf{Odróżnianie bliskich harmonicznie dźwięków:} W gęstych strukturach harmonicznych, rozróżnienie blisko brzmiących dźwięków może być trudne dla systemów ATM \cite{bib:benetos2013automatic}.
    
    \item \textbf{Interpretacja rytmiczna:} Precyzyjne odwzorowanie złożonych struktur rytmicznych, szczególnie w muzyce z rubato lub improwizowanej, stanowi wyzwanie dla obecnych systemów \cite{bib:cheng2016attack}.
    
    \item \textbf{Wykorzystanie pedałów:} Uwzględnienie u użytkowania pedałów fortepianowych, które wpływają na rezonans strun i brzmienie, wymaga zaawansowanych technik analizy akustycznej \cite{bib:cheng2016attack}.
\end{itemize}

Projekt TuneScript wpisuje się w ten kontekst, stawiając sobie za cel adresowanie wymienionych wyzwań poprzez integrację zaawansowanego modelu AI do transkrypcji z MP3 do MIDI z narzędziami do generowania partytur PDF. Takie kompleksowe podejście ma potencjał przezwyciężenia niektórych ograniczeń obecnych rozwiązań i dostarczenia użytkownikom narzędzia o wysokiej praktycznej wartości.

\section{Studia literaturowe}

W celu głębszego zrozumienia problematyki automatycznej transkrypcji muzyki fortepianowej oraz określenia aktualnego stanu wiedzy w tej dziedzinie, przeprowadzono szczegółowy przegląd literatury. Poniżej przedstawiono kluczowe prace i ich wkład w rozwój tej dziedziny:

\begin{enumerate}
    \item \textbf{Kong et al. (2020) \cite{bib:kong2020high}} w pracy "High-resolution Piano Transcription with Pedals by Regressing Onsets and Offsets Times" prezentują nowatorskie podejście do transkrypcji fortepianowej, które stanowi podstawę modelu wykorzystywanego w projekcie TuneScript. Autorzy proponują metodę wysokiej rozdzielczości czasowej do transkrypcji fortepianowej, która uwzględnia nie tylko detekcję nut, ale także użycie pedałów. Kluczową innowacją jest zastosowanie regresji do precyzyjnego określania czasów początku (onset) i końca (offset) każdej nuty, co pozwala na dokładniejsze odwzorowanie rytmu i artykulacji. Model ten wykazuje znaczącą poprawę w stosunku do wcześniejszych rozwiązań, szczególnie w zakresie precyzji czasowej i detekcji użycia pedałów.

    \item \textbf{Benetos et al. (2013) \cite{bib:benetos2013automatic}} w swojej pracy przeglądowej "Automatic music transcription: challenges and future directions" przedstawiają kompleksowe omówienie technik ATM. Autorzy omawiają główne wyzwania w dziedzinie, takie jak polifonia, różnorodność instrumentów i style wykonania. W kontekście muzyki fortepianowej, praca ta podkreśla złożoność problemu transkrypcji ze względu na szeroką gamę dynamiczną i bogactwo harmoniczne fortepianu.
    
    \item \textbf{Hawthorne et al. (2018) \cite{bib:hawthorne2018onsets}} w artykule "Onsets and frames: Dual-objective piano transcription" prezentują przełomowy model do transkrypcji fortepianowej. Ich podejście, nazwane "Onsets and Frames", wykorzystuje osobne sieci neuronowe do detekcji początków dźwięków i ich trwania. Model ten znacząco poprawił dokładność transkrypcji fortepianowej, osiągając wyniki bliskie ludzkiej dokładności w niektórych aspektach. Jest to jedno z najbardziej zaawansowanych rozwiązań w dziedzinie ATM dla fortepianu, które stanowiło inspirację dla późniejszych prac, w tym modelu Kong et al.
    
    \item \textbf{Kelz et al. (2016) \cite{bib:kelz2016potential}} w pracy "On the Potential of Simple Framewise Approaches to Piano Transcription" badają możliwości prostszych podejść do transkrypcji fortepianowej, opierających się na analizie pojedynczych ramek czasowych. Autorzy wykazują, że nawet stosunkowo proste modele mogą osiągać konkurencyjne wyniki, co ma istotne implikacje dla projektowania systemów działających w czasie rzeczywistym.
    
    \item \textbf{Cheng et al. (2016) \cite{bib:cheng2016attack}} w artykule "Attack/decay analysis of piano sounds using an exponential model" skupiają się na modelowaniu akustycznych właściwości dźwięków fortepianowych. Ich praca jest istotna dla zrozumienia, jak precyzyjnie analizować początkową fazę dźwięku (atak) oraz jego wygaszanie (decay), co jest kluczowe dla dokładnej transkrypcji dynamiki i artykulacji.
    
    \item \textbf{Jeong et al. (2017) \cite{bib:jeong2017deep}} w pracy "Deep Learning for Music" prezentują szerokie zastosowanie technik głębokiego uczenia w różnych zadaniach związanych z przetwarzaniem muzyki, w tym w transkrypcji fortepianowej. Autorzy omawiają, jak różne architektury sieci neuronowych mogą być wykorzystane do modelowania różnych aspektów muzyki, od wysokości dźwięku po ekspresję.
    
    \item \textbf{Sigtia et al. (2016) \cite{bib:sigtia2016end}} w pracy "An end-to-end neural network for polyphonic piano music transcription" prezentują kompleksowe podejście do transkrypcji muzyki fortepianowej wykorzystujące głębokie sieci neuronowe. Ich model jest w stanie bezpośrednio przetwarzać surowy sygnał audio na notację muzyczną, co stanowi znaczący krok w kierunku w pełni automatycznej transkrypcji.
\end{enumerate}

Powyższe prace stanowią fundament teoretyczny dla projektu TuneScript. Pokazują one ewolucję technik ATM dla fortepianu od tradycyjnych metod przetwarzania sygnałów do zaawansowanych modeli uczenia głębokiego. Jednocześnie, literatura ta wskazuje na wciąż istniejące wyzwania, szczególnie w zakresie transkrypcji muzyki polifonicznej, dokładnego odwzorowania dynamiki i ekspresji oraz generowania wysokiej jakości partytur.

Projekt TuneScript czerpie z tych badań, integrując najnowsze osiągnięcia w dziedzinie ATM z praktycznym aspektem generowania użytecznych partytur. Szczególnie istotna jest praca Kong et al. \cite{bib:kong2020high}, która stanowi podstawę modelu transkrypcji wykorzystywanego w systemie. Model ten, dzięki swojej wysokiej rozdzielczości czasowej i uwzględnieniu pedałów, pozwala na bardziej precyzyjną transkrypcję niż wcześniejsze rozwiązania.

Warto zauważyć, że mimo znaczącego postępu w dziedzinie ATM dla fortepianu, wciąż istnieje przestrzeń do dalszych ulepszeń, szczególnie w zakresie dokładności transkrypcji bardzo złożonych utworów polifonicznych oraz wiernego odwzorowania subtelności ekspresji muzycznej. TuneScript stara się adresować te wyzwania, oferując kompleksowe rozwiązanie łączące zaawansowaną transkrypcję z praktycznym zastosowaniem w postaci generowania partytur.
%%%%%%%%%%%%%%%%%%%%%%%%




% TODO
\chapter{Wymagania i narzędzia}
\label{ch:wymagania-i-narzedzia}

\section{Wymagania funkcjonalne}

System TuneScript powinien spełniać następujące wymagania funkcjonalne:

\begin{enumerate}
    \item \textbf{Uwierzytelnianie i autoryzacja użytkowników:}
    \begin{itemize}
        \item System musi zapewniać funkcjonalności rejestracji, logowania i wylogowywania użytkowników.
        \item Użytkownicy muszą mieć możliwość resetowania haseł za pomocą poczty elektronicznej.
    \end{itemize}
    
    \item \textbf{Przesyłanie plików:}
    \begin{itemize}
        \item System musi umożliwiać użytkownikom przesyłanie plików audio w formatach MP3 i WAV.
        \item System musi weryfikować format przesyłanych plików.
    \end{itemize}
    
    \item \textbf{Przetwarzanie transkrypcji:}
    \begin{itemize}
        \item System musi przetwarzać przesłane pliki audio w celu generowania transkrypcji fortepianowych.
        \item Transkrypcje powinny być zapisywane w formacie łatwym do odczytu i edycji.
    \end{itemize}
    
    \item \textbf{Zarządzanie transkrypcjami:}
    \begin{itemize}
        \item Użytkownicy muszą mieć możliwość przeglądania i usuwania swoich transkrypcji.
        \item Użytkownicy muszą mieć możliwość ustawiania widoczności swoich transkrypcji jako publicznych lub prywatnych.
    \end{itemize}
    
    \item \textbf{Funkcjonalność wyszukiwania:}
    \begin{itemize}
        \item System musi zapewniać podstawowe możliwości wyszukiwania transkrypcji na podstawie tytułu, kompozytora i widoczności.
    \end{itemize}
    
    \item \textbf{System oceniania:}
    \begin{itemize}
        \item Użytkownicy muszą mieć możliwość oceniania transkrypcji.
        \item System musi obliczać i wyświetlać średnią ocenę dla każdej transkrypcji.
    \end{itemize}
    
    \item \textbf{Zakładki:}
    \begin{itemize}
        \item Użytkownicy muszą mieć możliwość dodawania transkrypcji do ulubionych.
        \item Użytkownicy muszą mieć możliwość przeglądania i zarządzania swoimi ulubionymi transkrypcjami.
    \end{itemize}
\end{enumerate}

\section{Wymagania niefunkcjonalne}

System TuneScript powinien spełniać następujące wymagania niefunkcjonalne:

\begin{enumerate}
    \item \textbf{Wydajność:}
    \begin{itemize}
        \item System musi przetwarzać transkrypcje w rozsądnych ramach czasowych (np. mniej niż 5 minut dla 10-minutowego pliku audio).
    \end{itemize}
    
    \item \textbf{Użyteczność:}
    \begin{itemize}
        \item Interfejs użytkownika musi być intuicyjny i łatwy w nawigacji.
        \item System musi dostarczać użytkownikom jasne instrukcje i informacje zwrotne.
    \end{itemize}
    
    \item \textbf{Niezawodność:}
    \begin{itemize}
        \item System musi być niezawodny, z minimalnym czasem przestoju.
        \item System musi zapewniać solidne mechanizmy obsługi błędów i odzyskiwania.
    \end{itemize}
    
    \item \textbf{Bezpieczeństwo:}
    \begin{itemize}
        \item System musi chronić dane użytkowników za pomocą odpowiednich środków bezpieczeństwa.
        \item System musi wykorzystywać bezpieczne mechanizmy uwierzytelniania.
    \end{itemize}
    
    \item \textbf{Kompatybilność:}
    \begin{itemize}
        \item System musi być kompatybilny z głównymi przeglądarkami internetowymi (np. Chrome, Firefox).
    \end{itemize}
    
    \item \textbf{Łatwość utrzymania:}
    \begin{itemize}
        \item Kod źródłowy musi być dobrze udokumentowany i modularny, aby ułatwić konserwację i aktualizacje.
        \item System musi być łatwy do wdrożenia i skalowania.
    \end{itemize}
\end{enumerate}

\section{Przypadki użycia}

Poniżej przedstawiono diagram przypadków użycia dla systemu TuneScript w notacji UML:

% \begin{figure}[h]
%     \centering
%     \includegraphics[width=0.8\textwidth]{use_case_diagram.png}
%     \caption{Diagram przypadków użycia systemu TuneScript}
%     \label{fig:use_case_diagram}
% \end{figure}

Główne przypadki użycia obejmują:
\begin{itemize}
    \item Rejestracja użytkownika
    \item Logowanie użytkownika
    \item Przesyłanie pliku audio
    \item Generowanie transkrypcji
    \item Edycja transkrypcji
    \item Wyszukiwanie transkrypcji
    \item Ocenianie transkrypcji
    \item Dodawanie transkrypcji do ulubionych
    \item Zarządzanie ulubionymi transkrypcjami
\end{itemize}

\section{Opis narzędzi}

Do realizacji projektu TuneScript wykorzystano następujące narzędzia i technologie:

\begin{itemize}
    \item \textbf{Backend:}
    \begin{itemize}
        \item Python
        \item Django: Framework webowy
        \item Graphene: Biblioteka do implementacji API GraphQL
        \item Celery: System kolejkowania zadań
        \item Redis: Broker wiadomości dla Celery
        \item PyTorch: Biblioteka do uczenia maszynowego
        \item Librosa: Biblioteka do analizy audio
        \item Flask: Lekki framework webowy
        \item Transformers: Biblioteka do przetwarzania języka naturalnego
        \item PostgreSQL: System zarządzania bazą danych
        \item MuseScore: Generowanie partytur w formacie PDF
    \end{itemize}
    
    \item \textbf{Frontend:}
    \begin{itemize}
        \item Next.js: Framework React do renderowania po stronie serwera
        \item TypeScript: Typowany nadzbiór JavaScript
        \item Tailwind CSS: Framework CSS utility-first
        \item Apollo: Klient GraphQL
    \end{itemize}
    
    \item \textbf{Narzędzia deweloperskie:}
    \begin{itemize}
        \item Poetry: Narzędzie do zarządzania zależnościami Python
        \item Docker: Platforma do konteneryzacji
        \item Docker Compose: Narzędzie do definiowania i uruchamiania aplikacji wielokontenerowych
        \item ESLint: Narzędzie do analizy statycznej kodu JavaScript
        \item Black: Formatter kodu Python
        \item isort: Narzędzie do sortowania importów w Pythonie
        \item Flake8: Linter kodu Python
    \end{itemize}
    
    \item \textbf{Narzędzia testowe:}
    \begin{itemize}
        \item Factory Boy: Biblioteka do generowania obiektów testowych w Pythonie
        \item pytest: Framework testowy dla Pythona
        \item Jest: Framework testowy dla JavaScript
    \end{itemize}
\end{itemize}

\section{Metodyka pracy nad projektowaniem i implementacją}

Proces projektowania i implementacji systemu TuneScript był realizowany w sposób iteracyjny, z naciskiem na stopniowe rozwijanie funkcjonalności. Główne etapy pracy obejmowały:

\begin{enumerate}
    \item \textbf{Analiza wymagań}: Szczegółowe określenie wymagań funkcjonalnych i niefunkcjonalnych na podstawie analizy potrzeb użytkowników i badania istniejących rozwiązań.
    
    \item \textbf{Projektowanie architektury}: Opracowanie ogólnej architektury systemu, z uwzględnieniem modułów odpowiedzialnych za transkrypcję, API backendu i interfejs użytkownika.
    
    \item \textbf{Implementacja i szkolenie modelu transkrypcji}: 
    \begin{itemize}
        \item Adaptacja modelu opartego na pracy Kong et al. \cite{bib:kong2020high} do potrzeb systemu TuneScript.
        \item Implementacja modelu z wykorzystaniem PyTorch i Librosa.
        \item Przeprowadzenie procesu szkolenia i optymalizacji modelu.
    \end{itemize}
    
    \item \textbf{Iteracyjny rozwój funkcjonalności}:
    Dla każdej funkcjonalności systemu proces obejmował:
    \begin{itemize}
        \item Implementację logiki backendu z wykorzystaniem Django i Graphene.
        \item Rozwój odpowiadającego interfejsu użytkownika w Next.js i TypeScript.
        \item Integrację backendu z frontendem za pomocą Apollo Client.
        \item Pisanie testów jednostkowych i integracyjnych dla zaimplementowanej funkcjonalności.
    \end{itemize}
    
    \item \textbf{Testowanie i debugowanie}: Po implementacji każdej funkcjonalności przeprowadzano testy, wykorzystując pytest dla backendu i Jest dla frontendu. W razie potrzeby przeprowadzano debugowanie i optymalizację kodu.
    
    \item \textbf{Integracja komponentów}: Stopniowe łączenie poszczególnych modułów systemu, w tym integracja modelu transkrypcji z API backendu i implementacja systemu kolejkowania zadań z użyciem Celery i Redis.
    
    \item \textbf{Optymalizacja wydajności}: Analiza wydajności systemu i optymalizacja krytycznych ścieżek, szczególnie w procesie transkrypcji i generowania partytur.
    
    \item \textbf{Dokumentacja}: Tworzenie dokumentacji technicznej równolegle z rozwojem systemu, obejmującej zarówno kod źródłowy, jak i instrukcje użytkownika.
    
    \item \textbf{Przygotowanie do wdrożenia}: Konfiguracja środowiska produkcyjnego z wykorzystaniem Docker i Docker Compose, zapewniająca łatwe wdrożenie i skalowalność systemu.
\end{enumerate}

Podczas całego procesu rozwoju projektu stosowano następujące praktyki:

\begin{itemize}
    \item \textbf{Kontrola wersji}: Wykorzystanie systemu Git do śledzenia zmian w kodzie i zarządzania wersjami.
    
    \item \textbf{Ciągła integracja}: Regularne integrowanie nowych funkcjonalności z główną gałęzią projektu, aby wcześnie wykrywać i rozwiązywać konflikty.
    
    \item \textbf{Przeglądy kodu}: Samodzielne przeglądy implementowanych funkcjonalności w celu zapewnienia jakości i spójności kodu.
    
    \item \textbf{Automatyzacja}: Wykorzystanie narzędzi takich jak Black, isort i Flake8 dla Pythona oraz ESLint dla JavaScript/TypeScript do automatycznego formatowania i lintowania kodu.
    
    \item \textbf{Zarządzanie zależnościami}: Użycie Poetry do zarządzania zależnościami Python, zapewniając reprodukowalność środowiska deweloperskiego.
\end{itemize}

Takie podejście pozwoliło na elastyczne rozwijanie projektu, umożliwiając skupienie się na jednej funkcjonalności na raz i zapewniając solidne podstawy dla każdego kolejnego etapu rozwoju systemu TuneScript.

% TODO
\chapter{Specyfikacja zewnętrzna}
\label{ch:04}

\section{Wymagania sprzętowe i programowe}

\subsection{Wymagania sprzętowe}
\begin{itemize}
    \item Serwer: 
    \begin{itemize}
        \item Procesor: minimum 4 rdzenie, zalecane 8 rdzeni
        \item RAM: minimum 8 GB, zalecane 16 GB
        \item Dysk: minimum 100 GB przestrzeni dyskowej SSD
    \end{itemize}
    \item Klient:
    \begin{itemize}
        \item Dowolne urządzenie z przeglądarką internetową i dostępem do internetu
    \end{itemize}
\end{itemize}

\subsection{Wymagania programowe}
\begin{itemize}
    \item Serwer:
    \begin{itemize}
        \item System operacyjny: Linux (zalecany Ubuntu 20.04 LTS lub nowszy)
        \item Docker w wersji 20.10 lub nowszej
        \item Docker Compose w wersji 1.29 lub nowszej
    \end{itemize}
    \item Klient:
    \begin{itemize}
        \item Nowoczesna przeglądarka internetowa (Chrome 90+, Firefox 88+, Safari 14+, Edge 90+)
    \end{itemize}
\end{itemize}

\section{Sposób instalacji}

Instalacja systemu TuneScript jest prosta i wymaga jedynie kilku kroków:

\begin{enumerate}
    \item Sklonuj repozytorium projektu:
    \begin{verbatim}
    git clone https://github.com/mj300405/TuneScript.git
    \end{verbatim}
    
    \item Przejdź do katalogu narzędzi deweloperskich projektu:
    \begin{verbatim}
    cd TuneScript/devtools
    \end{verbatim}
    
    \item Uruchom Docker Compose, aby zbudować i uruchomić kontenery (upewnij się, że masz zainstalowany Docker):
    \begin{verbatim}
    docker-compose up --build
    \end{verbatim}
    
    \item Po zakończeniu procesu budowania i uruchamiania, aplikacja będzie dostępna w przeglądarce pod adresem:
    \begin{verbatim}
    http://localhost:3000
    \end{verbatim}
\end{enumerate}

\textbf{Uwaga:} Przed uruchomieniem Docker Compose należy upewnić się, że Docker jest zainstalowany na systemie. Jeśli nie, można pobrać go i zainstalować ze strony oficjalnej: \url{https://www.docker.com/get-started}.

\section{Kategorie użytkowników}

\begin{itemize}
    \item \textbf{Niezarejestrowani użytkownicy}: Mogą przeglądać publiczne transkrypcje.
    \item \textbf{Zarejestrowani użytkownicy}: Mogą przesyłać pliki audio, generować transkrypcje, zarządzać własnymi transkrypcjami, oceniać i dodawać do ulubionych transkrypcje innych użytkowników.
    \item \textbf{Administratorzy}: Mają pełny dostęp do systemu, w tym do panelu administracyjnego.
\end{itemize}

\section{Sposób obsługi}

\begin{enumerate}
    \item \textbf{Rejestracja i logowanie}:
    \begin{itemize}
        \item Użytkownik może zarejestrować się, podając nazwę użytkownika, adres e-mail i hasło.
        \item Następnie otrzymuje e-mail z linkiem potwierdzającym.
        \item Po potiwerdzeniu, użytkownik może zalogować się do systemu.
    \end{itemize}
    
    \item \textbf{Przesyłanie pliku audio}:
    \begin{itemize}
        \item Zalogowany użytkownik może przesłać plik audio w formacie MP3 lub WAV.
        \item System weryfikuje format i rozmiar pliku.
    \end{itemize}
    
    \item \textbf{Generowanie transkrypcji}:
    \begin{itemize}
        \item Po przesłaniu pliku, system automatycznie rozpoczyna proces transkrypcji.
        \item Użytkownik jest informowany o postępie procesu.
    \end{itemize}
    
    \item \textbf{Zarządzanie transkrypcjami}:
    \begin{itemize}
        \item Użytkownik może przeglądać, edytować i usuwać swoje transkrypcje.
        \item Możliwe jest ustawienie widoczności transkrypcji jako publicznej lub prywatnej.
    \end{itemize}
    
    \item \textbf{Wyszukiwanie i przeglądanie transkrypcji}:
    \begin{itemize}
        \item Użytkownicy mogą wyszukiwać transkrypcje po tytule, kompozytorze i widoczności.
        \item Możliwe jest przeglądanie publicznych transkrypcji innych użytkowników.
    \end{itemize}
    
    \item \textbf{Ocenianie i dodawanie do ulubionych}:
    \begin{itemize}
        \item Zalogowani użytkownicy mogą oceniać transkrypcje.
        \item Możliwe jest dodawanie transkrypcji do listy ulubionych.
    \end{itemize}
\end{enumerate}

\section{Administracja systemem}

Administratorzy mają dostęp do panelu administracyjnego Django, który umożliwia:
\begin{itemize}
    \item Zarządzanie kontami użytkowników
    \item Przeglądanie i moderowanie transkrypcji
    \item Monitorowanie aktywności systemu
    \item Konfigurację ustawień systemu
\end{itemize}

\section{Kwestie bezpieczeństwa}

\begin{itemize}
    \item Wszystkie hasła są przechowywane w formie zaszyfrowanej.
    \item Komunikacja między klientem a serwerem jest szyfrowana za pomocą HTTPS.
    \item System wykorzystuje tokeny JWT do uwierzytelniania użytkowników.
    \item Implementacja mechanizmów ochrony przed atakami CSRF i XSS.
    \item Regularne aktualizacje wszystkich wykorzystywanych bibliotek i zależności.
\end{itemize}

\section{Przykład działania}

Poniżej przedstawiono przykładowy scenariusz korzystania z systemu TuneScript:

\begin{enumerate}
    \item Użytkownik rejestruje się w systemie.
    \item Po zalogowaniu, użytkownik przesyła plik audio zawierający nagranie fortepianowe.
    \item System przetwarza plik i generuje transkrypcję.
    \item Użytkownik może przeglądać, edytować i pobierać wygenerowaną transkrypcję.
    \item Użytkownik ustawia transkrypcję jako publiczną.
    \item Inni użytkownicy mogą przeglądać, oceniać i dodawać transkrypcję do ulubionych.
\end{enumerate}

\section{Scenariusze korzystania z systemu}

% \begin{figure}[h]
%     \centering
%     \includegraphics[width=0.8\textwidth]{registration_screen.png}
%     \caption{Ekran rejestracji użytkownika}
%     \label{fig:registration_screen}
% \end{figure}

% \begin{figure}[h]
%     \centering
%     \includegraphics[width=0.8\textwidth]{upload_screen.png}
%     \caption{Ekran przesyłania pliku audio}
%     \label{fig:upload_screen}
% \end{figure}

% \begin{figure}[h]
%     \centering
%     \includegraphics[width=0.8\textwidth]{transcription_view.png}
%     \caption{Widok wygenerowanej transkrypcji}
%     \label{fig:transcription_view}
% \end{figure}

% \begin{figure}[h]
%     \centering
%     \includegraphics[width=0.8\textwidth]{search_results.png}
%     \caption{Wyniki wyszukiwania transkrypcji}
%     \label{fig:search_results}
% \end{figure}



% TODO
\chapter{[Właściwy dla kierunku -- np. Specyfikacja wewnętrzna]}
\label{ch:05}


Jeśli „Specyfikacja wewnętrzna”:
\begin{itemize}
\item przedstawienie idei
\item architektura systemu
\item opis struktur danych (i organizacji baz danych)
\item komponenty, moduły, biblioteki, przegląd ważniejszych klas (jeśli występują)
\item przegląd ważniejszych algorytmów (jeśli występują)
\item szczegóły implementacji wybranych fragmentów, zastosowane wzorce projektowe
\item diagramy UML
\end{itemize}

% % % % % % % % % % % % % % % % % % % % % % % % % % % % % % % % % % % 
% Pakiet minted wymaga importu: \usepackage{minted}                 %
% i specjalnego kompilowania:                                       %
% pdflatex -shell-escape main                                       %
% % % % % % % % % % % % % % % % % % % % % % % % % % % % % % % % % % % 


Krótka wstawka kodu w linii tekstu jest możliwa, np.  \lstinline|int a;| (biblioteka \texttt{listings})% lub  \mintinline{C++}|int a;| (biblioteka \texttt{minted})
. 
Dłuższe fragmenty lepiej jest umieszczać jako rysunek, np. kod na rys \ref{fig:pseudokod:listings}% i rys. \ref{fig:pseudokod:minted}
, a naprawdę długie fragmenty – w załączniku.


\begin{figure}
\centering
\begin{lstlisting}
class test : public basic
{
    public:
      test (int a);
      friend std::ostream operator<<(std::ostream & s, 
                                     const test & t);
    protected:
      int _a;  
      
};
\end{lstlisting}
\caption{Pseudokod w \texttt{listings}.}
\label{fig:pseudokod:listings}
\end{figure}

%\begin{figure}
%\centering
%\begin{minted}[linenos,frame=lines]{c++}
%class test : public basic
%{
%    public:
%      test (int a);
%      friend std::ostream operator<<(std::ostream & s, 
%                                     const test & t);
%    protected:
%      int _a;  
%      
%};
%\end{minted}
%\caption{Pseudokod w \texttt{minted}.}
%\label{fig:pseudokod:minted}
%\end{figure}




% TODO
\chapter{Weryfikacja i walidacja}
\label{ch:06}
\begin{itemize}
\item sposób testowania w ramach pracy (np. odniesienie do modelu V)
\item organizacja eksperymentów
\item przypadki testowe zakres testowania (pełny/niepełny)
\item wykryte i usunięte błędy
\item opcjonalnie wyniki badań eksperymentalnych
\end{itemize}

\begin{table}
\centering
\caption{Nagłówek tabeli jest nad tabelą.}
\label{id:tab:wyniki}
\begin{tabular}{rrrrrrrr}
\toprule
	         &                                     \multicolumn{7}{c}{metoda}                                      \\
	         \cmidrule{2-8}
	         &         &         &        \multicolumn{3}{c}{alg. 3}        & \multicolumn{2}{c}{alg. 4, $\gamma = 2$} \\
	         \cmidrule(r){4-6}\cmidrule(r){7-8}
	$\zeta$ &     alg. 1 &   alg. 2 & $\alpha= 1.5$ & $\alpha= 2$ & $\alpha= 3$ &   $\beta = 0.1$  &   $\beta = -0.1$ \\
\midrule
	       0 &  8.3250 & 1.45305 &       7.5791 &    14.8517 &    20.0028 & 1.16396 &                       1.1365 \\
	       5 &  0.6111 & 2.27126 &       6.9952 &    13.8560 &    18.6064 & 1.18659 &                       1.1630 \\
	      10 & 11.6126 & 2.69218 &       6.2520 &    12.5202 &    16.8278 & 1.23180 &                       1.2045 \\
	      15 &  0.5665 & 2.95046 &       5.7753 &    11.4588 &    15.4837 & 1.25131 &                       1.2614 \\
	      20 & 15.8728 & 3.07225 &       5.3071 &    10.3935 &    13.8738 & 1.25307 &                       1.2217 \\
	      25 &  0.9791 & 3.19034 &       5.4575 &     9.9533 &    13.0721 & 1.27104 &                       1.2640 \\
	      30 &  2.0228 & 3.27474 &       5.7461 &     9.7164 &    12.2637 & 1.33404 &                       1.3209 \\
	      35 & 13.4210 & 3.36086 &       6.6735 &    10.0442 &    12.0270 & 1.35385 &                       1.3059 \\
	      40 & 13.2226 & 3.36420 &       7.7248 &    10.4495 &    12.0379 & 1.34919 &                       1.2768 \\
	      45 & 12.8445 & 3.47436 &       8.5539 &    10.8552 &    12.2773 & 1.42303 &                       1.4362 \\
	      50 & 12.9245 & 3.58228 &       9.2702 &    11.2183 &    12.3990 & 1.40922 &                       1.3724 \\
\bottomrule
\end{tabular}
\end{table}  



% TODO
\chapter{Podsumowanie i wnioski}
\begin{itemize}
\item uzyskane wyniki w świetle postawionych celów i zdefiniowanych wyżej wymagań
\item kierunki ewentualnych danych prac (rozbudowa funkcjonalna …)
\item problemy napotkane w trakcie pracy
\end{itemize}



\backmatter

%\bibliographystyle{plplain}  % bibtex
%\bibliography{biblio} % bibtex
\printbibliography           % biblatex
\addcontentsline{toc}{chapter}{Bibliografia}

\begin{appendices}

% TODO
\chapter{Spis skrótów i symboli}

\begin{itemize}
\item[DNA] kwas deoksyrybonukleinowy (ang. \english{deoxyribonucleic acid})
\item[MVC] model -- widok -- kontroler (ang. \english{model--view--controller}) 
\item[$N$] liczebność zbioru danych
\item[$\mu$] stopnień przyleżności do zbioru
\item[$\mathbb{E}$] zbiór krawędzi grafu
\item[$\mathcal{L}$] transformata Laplace'a 
\end{itemize}


% TODO
\chapter{Źródła}

Jeżeli w pracy konieczne jest umieszczenie długich fragmentów kodu źródłowego, należy je przenieść w to miejsce.

\begin{lstlisting}
if (_nClusters < 1)
	throw std::string ("unknown number of clusters");
if (_nIterations < 1 and _epsilon < 0)
	throw std::string ("You should set a maximal number of iteration or minimal difference -- epsilon.");
if (_nIterations > 0 and _epsilon > 0)
	throw std::string ("Both number of iterations and minimal epsilon set -- you should set either number of iterations or minimal epsilon.");
\end{lstlisting}


% % % % % % % % % % % % % % % % % % % % % % % % % % % % % % % % % % % 
% Pakiet minted wymaga odkomentowania w pliku config/settings.tex   %
% importu pakietu minted: \usepackage{minted}                       %
% i specjalnego kompilowania:                                       %
% pdflatex -shell-escape praca                                      %
% % % % % % % % % % % % % % % % % % % % % % % % % % % % % % % % % % % 

%\begin{minted}[linenos,breaklines,frame=lines]{c++}
%if (_nClusters < 1)
%   throw std::string ("unknown number of clusters");
%if (_nIterations < 1 and _epsilon < 0)
%   throw std::string ("You should set a maximal number of iteration or minimal difference -- epsilon.");
%if (_nIterations > 0 and _epsilon > 0)
%   throw std::string ("Both number of iterations and minimal epsilon set -- you should set either number of iterations or minimal epsilon.");
%\end{minted}


% TODO
\chapter{Lista dodatkowych plików, uzupełniających tekst pracy} 


W systemie do pracy dołączono dodatkowe pliki zawierające:
\begin{itemize}
\item źródła programu,
\item dane testowe,
\item film pokazujący działanie opracowanego oprogramowania lub zaprojektowanego i~wykonanego urządzenia,
\item itp.
\end{itemize}


\listoffigures
\addcontentsline{toc}{chapter}{Spis rysunków}
\listoftables
\addcontentsline{toc}{chapter}{Spis tabel}

\end{appendices}

\end{document}


%% Finis coronat opus.

